\documentclass[11pt]{amsart}
\usepackage[T1]{fontenc}
\usepackage{lmodern,microtype,amsmath,amssymb,amsthm,mathtools,enumitem,booktabs}
\usepackage[margin=1.1in]{geometry}
\usepackage[hidelinks]{hyperref}
\usepackage{xurl}

\newtheorem{theorem}{Theorem}[section]
\newtheorem{lemma}[theorem]{Lemma}
\newtheorem{corollary}[theorem]{Corollary}
\newtheorem{proposition}[theorem]{Proposition}
\theoremstyle{definition}
\newtheorem{definition}[theorem]{Definition}
\theoremstyle{remark}
\newtheorem{remark}[theorem]{Remark}
\numberwithin{equation}{section}

\DeclareMathOperator{\diam}{diam}
\DeclareMathOperator{\dist}{dist}
\DeclareMathOperator{\tr}{tr}

\hypersetup{pdftitle={The largest Kemeny constant and the largest relaxation time of a graph},pdfauthor={Yitzchak Shmalo}}

\title[The largest Kemeny constant of a graph]{The largest Kemeny constant and the largest relaxation time of a graph}
\author{Yitzchak Shmalo}
\address{Einstein Institute of Mathematics, The Hebrew University of Jerusalem, Jerusalem, Israel}
\email{yitzchak.shmalo@gmail.com}
\date{6 November 2026}
\subjclass[2020]{60J10, 05C81, 05C50, 05C35}
\keywords{random walk on a graph, Kemeny's constant, relaxation time, spectral gap, effective resistance, barbell graph}

\begin{document}
\begin{abstract}
For the simple random walk on a connected graph with $n$ vertices, let $\tau_0$ be Kemeny's constant and $\tau_2$ the relaxation time. Aldous and Fill conjectured that the maxima of both quantities over all connected $n$-vertex graphs are asymptotic to $n^3/54$. We prove this. Every connected graph on $n$ vertices satisfies $\tau_2\le\tau_0\le n^3/54+\tfrac{37}{8}n^{5/2}$, and the barbell made of two cliques on about $n/3$ vertices joined by a path on about $n/3$ vertices has $\tau_2\ge(n-2)(n-5)(n+3)/54$. The upper bound comes from the electrical formula for Kemeny's constant: the vertices near high-degree vertices are grouped into components of small total diameter, the rest of the graph is compared with a tree, and each edge of the tree is shown to contribute at most $s^2/8$ plus a lower-order term, where $s$ is the number of grouped vertices.
\end{abstract}
\maketitle

\section{Introduction}\label{sec:intro}

Let $G=(V,E)$ be a finite connected simple graph with $n\ge2$ vertices, degrees $d_v$, and total degree $W=\sum_v d_v=2|E|$. The simple random walk on $G$ has transition matrix $P(u,v)=1/d_u$ for $uv\in E$ and stationary distribution $\pi_v=d_v/W$. We run it in continuous time with generator $P-I$, so that it jumps at rate one; expected hitting times are the same as for the discrete-time walk. Let $T_v$ be the hitting time of $v$, with $T_v=0$ for a walk started at $v$, and let
\[
0=\mu_1<\mu_2\le\dots\le\mu_n
\]
be the eigenvalues of $I-P$. Following Aldous and Fill~\cite{AF}, put
\[
\tau_0=\sum_{v\in V}\pi_v E_u[T_v],\qquad \tau_2=\frac1{\mu_2}.
\]
The sum $\tau_0$ does not depend on $u$; it is Kemeny's constant of the walk~\cite{KS}. The relaxation time $\tau_2$ is the inverse spectral gap. Some authors define Kemeny's constant with an extra term $1$, or for the lazy walk $(I+P)/2$, which doubles both $\tau_0$ and $\tau_2$; the constant $1/54$ below refers to the normalisation just given.

Among the extremal questions collected by Aldous and Fill, the largest hitting time and the largest cover time of an $n$-vertex graph are of order $n^3$, with constant $4/27$, attained asymptotically by lollipop graphs (Brightwell and Winkler~\cite{BW}; Feige~\cite{Fe}). For $\tau_0$ and $\tau_2$ the natural candidates are barbells: two cliques joined by a long path. Aldous and Fill~\cite[Open Problem~6.13]{AF} conjecture that
\[
\max_G\tau_0(G)\sim\frac{n^3}{54}\quad\text{and}\quad\max_G\tau_2(G)\sim\frac{n^3}{54},
\]
the maxima being over connected graphs with $n$ vertices. Breen, Butler, Day, DeArmond, Lorenzen, Qian and Riesen~\cite{BBD} computed Kemeny's constant of barbell-type graphs and showed that it is of order $n^3$, the largest possible order. Faught, Kempton and Knudson~\cite{FKK} obtained Kemeny's constant of a barbell from a formula for graphs with a cut vertex and showed that the best barbell has cliques on about $n/3$ vertices each and Kemeny's constant $n^3/54+O(n^2)$; they note that barbells are believed, on the basis of computation, to maximise Kemeny's constant. The matching upper bound is the content of the present paper.

\begin{theorem}\label{thm:main}
Let $G$ be a connected simple graph on $n\ge2$ vertices. Then
\begin{equation}\label{eq:upper}
\tau_2(G)\le\tau_0(G)\le\frac{n^3}{54}+\frac{37}{8}\,n^{5/2}.
\end{equation}
For $n\ge6$ let $a=\lfloor n/3\rfloor$ and $\ell=n-2a+1$, and let $B_n$ be the barbell consisting of two disjoint copies of $K_a$ whose chosen vertices are joined by a path with $\ell$ edges. Then $B_n$ has $n$ vertices and
\begin{equation}\label{eq:lower}
\tau_0(B_n)\ge\tau_2(B_n)\ge\frac{3a(a-1)\ell+\ell^2+2}{6}\ge\frac{(n-2)(n-5)(n+3)}{54}.
\end{equation}
\end{theorem}

\begin{corollary}\label{cor:asym}
As $n\to\infty$,
\[
\max_G\tau_0(G)=\frac{n^3}{54}+O(n^{5/2}),\qquad \max_G\tau_2(G)=\frac{n^3}{54}+O(n^{5/2}),
\]
the maxima being over connected simple graphs on $n$ vertices, and the barbells $B_n$ attain the leading constant for both quantities.
\end{corollary}

The constant has a simple explanation, and the proof follows it. In a barbell with cliques of size $a$ and a path of $\ell$ edges, each path edge separates two halves of total degree about $a^2$ each, and the electrical formula~\eqref{eq:resistance} below charges each such edge about $a^2\cdot a^2/(2a^2)=a^2/2$. The total, $a^2\ell/2$ with $2a+\ell\approx n$, is largest at $a=\ell=n/3$, where it equals $n^3/54$. For a general graph we single out the set $S$ of vertices within distance one of a vertex of degree at least $\sqrt n$. Its components have small total diameter, every other vertex has small degree, and after the components are contracted the graph is compared with a tree. A tree edge separates the $s=|S|$ grouped vertices into two parts of sizes $x$ and $y$, and its contribution is at most $x^2y^2/(x^2+y^2)\le s^2/8$, up to a term of order $n^{3/2}$. There are at most $n-s$ such edges, up to lower-order terms, and $(n-s)s^2/8\le n^3/54$.

The theorem leaves open the exact maximiser for each $n$ and the true size of the error term. The best barbell has $\tau_0=n^3/54+O(n^2)$ by~\cite{FKK}, which suggests that the error term in~\eqref{eq:upper} can be reduced to $O(n^2)$.

Section~\ref{sec:prelim} recalls the spectral and electrical formulas for $\tau_0$. Section~\ref{sec:upper} proves the upper bound~\eqref{eq:upper}, Section~\ref{sec:lower} the lower bound~\eqref{eq:lower}, and Section~\ref{sec:numerics} reports a numerical check.

\section{Spectral and electrical formulas}\label{sec:prelim}

Let $L=\operatorname{diag}(d)-A$ be the combinatorial Laplacian of $G$, where $A$ is the adjacency matrix. The matrix $I-P=\operatorname{diag}(d)^{-1}L$ is similar to the symmetric matrix $\operatorname{diag}(d)^{-1/2}L\operatorname{diag}(d)^{-1/2}$, and since $f^{\mathsf T}Lf=\sum_{uv\in E}(f_u-f_v)^2$ and $G$ is connected, its eigenvalue $0$ is simple. The variational characterisation gives
\begin{equation}\label{eq:rayleigh}
\mu_2=\min\Bigl\{\frac{\sum_{uv\in E}(f_u-f_v)^2}{\sum_v d_vf_v^2}:\ f\ne0,\ \sum_v d_vf_v=0\Bigr\}.
\end{equation}

\begin{lemma}\label{lem:formulas}
For the simple random walk on a connected graph,
\begin{equation}\label{eq:eigentime}
\tau_0=\sum_{i=2}^n\frac1{\mu_i},
\end{equation}
and in particular $\tau_2\le\tau_0$. Moreover
\begin{equation}\label{eq:resistance}
\tau_0=\frac1W\sum_{\{u,v\}\subset V}d_ud_vR_{uv},
\end{equation}
where the sum is over unordered pairs of distinct vertices and $R_{uv}$ is the effective resistance between $u$ and $v$ when every edge has unit resistance. Finally $R_{uv}\le\dist_G(u,v)$.
\end{lemma}

\begin{proof}
Identity~\eqref{eq:eigentime} is the eigentime identity~\cite[Chapter~3]{AF}; it also follows from the fundamental matrix $Z=(I-P+\mathbf 1\pi^{\mathsf T})^{-1}$, which has eigenvalues $1$ and $1/\mu_i$ ($i\ge2$) and satisfies $E_u[T_v]=(Z_{vv}-Z_{uv})/\pi_v$~\cite{KS}, so that $\tau_0=\tr Z-\sum_vZ_{uv}=\tr Z-1$. Since every term of~\eqref{eq:eigentime} is positive, $\tau_2=1/\mu_2\le\tau_0$.

For~\eqref{eq:resistance}, average the definition of $\tau_0$ over the starting vertex $u$ with respect to $\pi$:
\[
\tau_0=\sum_{u,v}\pi_u\pi_vE_u[T_v]=\frac12\sum_{u,v}\pi_u\pi_v\bigl(E_u[T_v]+E_v[T_u]\bigr).
\]
By the commute time identity~\cite{CRRST}, $E_u[T_v]+E_v[T_u]=WR_{uv}$. Hence $\tau_0=\frac W2\sum_{u,v}\pi_u\pi_vR_{uv}=\frac1{2W}\sum_{u,v}d_ud_vR_{uv}$, which is~\eqref{eq:resistance}.

By Thomson's principle, $R_{uv}$ is the minimum energy $\sum_{e}j_e^2$ of a unit flow $j$ from $u$ to $v$. The unit flow along a shortest path from $u$ to $v$ has energy $\dist_G(u,v)$, so $R_{uv}\le\dist_G(u,v)$.
\end{proof}

\section{The upper bound}\label{sec:upper}

Throughout this section $G$ is a connected simple graph on $n$ vertices and $D>0$ is a threshold. Let
\[
H=\{v\in V:d_v\ge D\},\qquad S=\bigcup_{h\in H}N[h],
\]
where $N[h]$ is the closed neighbourhood of $h$. Let $C_1,\dots,C_k$ be the vertex sets of the components of the induced graph $G[S]$, let $s_i=|C_i|$, $s=|S|=\sum_is_i$, and let $\delta_i$ be the diameter of the induced graph $G[C_i]$. Put $\delta=\sum_i\delta_i$. If $H$ is empty, then $S$ is empty and $s=k=\delta=0$.

\begin{lemma}\label{lem:components}
We have $k\le s/D$ and $\delta\le5s/D$.
\end{lemma}

\begin{proof}
Each $C_i$ contains a vertex $h\in H$ together with all of $N[h]$, because $N[h]\subset S$ is connected. Hence $s_i\ge|N[h]|>D$ and $k\le s/D$.

Fix $i$ and a shortest path $v_0,v_1,\dots,v_{\delta_i}$ in $G[C_i]$. For each index $j\in\{0,5,10,\dots\}$ with $j\le\delta_i$, choose $h_j\in H$ with $v_j\in N[h_j]$; such a vertex exists because $v_j\in S$, and then $N[h_j]\subset C_i$. If $N[h_j]$ and $N[h_m]$ met for two distinct selected indices $j<m$, the walk from $v_j$ through $h_j$, a common vertex of the two neighbourhoods and $h_m$ to $v_m$ would have length at most four and would stay in $C_i$, while $\dist_{G[C_i]}(v_j,v_m)=m-j\ge5$. So the sets $N[h_j]$ are pairwise disjoint subsets of $C_i$, each with more than $D$ elements, and
\[
\Bigl(\Bigl\lfloor\frac{\delta_i}5\Bigr\rfloor+1\Bigr)D\le s_i .
\]
Since $\delta_i/5<\lfloor\delta_i/5\rfloor+1$, this gives $\delta_i\le5s_i/D$. Summing over $i$ gives $\delta\le5s/D$.
\end{proof}

Contract each $C_i$ to a single vertex and keep every vertex outside $S$. The resulting graph $Q$ is connected and has $k+(n-s)$ vertices; let $q:V\to V(Q)$ be the quotient map. Fix a spanning tree $T$ of $Q$. It has
\begin{equation}\label{eq:edges}
|E(T)|=k+n-s-1\le n-s+\frac sD
\end{equation}
edges, and also $|E(T)|\le n-1$. For a vertex $x$ of $Q$ let $w_x=\sum_{v\in q^{-1}(x)}d_v$, so that $\sum_xw_x=W$. For an edge $e$ of $T$, let $W_L(e)$ and $W_R(e)$ be the sums of $w_x$ over the two components of $T-e$.

\begin{lemma}\label{lem:tree}
We have
\[
\tau_0\le\frac1W\sum_{e\in E(T)}W_L(e)W_R(e)+\frac{W\delta}2 .
\]
\end{lemma}

\begin{proof}
Let $u,v\in V$ and let $x_0=q(u),x_1,\dots,x_m=q(v)$ be the path from $q(u)$ to $q(v)$ in $T$. Each edge $x_{r-1}x_r$ of $T$ comes from an edge of $G$ between the parts $q^{-1}(x_{r-1})$ and $q^{-1}(x_r)$. Joining these $m$ edges by paths inside the contracted components, and adding a path inside the component of $u$ (or of $v$) when $u$ (or $v$) lies in $S$, gives a walk from $u$ to $v$ in $G$. The tree path visits each contracted component at most once, and inside $C_i$ the walk uses a path of length at most $\delta_i$. Therefore
\[
\dist_G(u,v)\le\dist_T(q(u),q(v))+\delta .
\]
By Lemma~\ref{lem:formulas},
\[
\tau_0\le\frac1W\sum_{\{u,v\}}d_ud_v\dist_T(q(u),q(v))+\frac\delta W\sum_{\{u,v\}}d_ud_v .
\]
The second term is at most $\delta W^2/(2W)=W\delta/2$. In the first, pairs inside one part contribute nothing, and grouping the rest by parts gives $\sum_{\{x,y\}}w_xw_y\dist_T(x,y)$. An edge $e$ of $T$ lies on the tree path between $x$ and $y$ exactly when it separates them, so this sum equals $\sum_eW_L(e)W_R(e)$.
\end{proof}

\begin{lemma}\label{lem:cut}
Every edge $e$ of $T$ satisfies
\[
\frac{W_L(e)W_R(e)}W\le\frac{s^2}8+2nD .
\]
\end{lemma}

\begin{proof}
For the vertex of $Q$ obtained from $C_i$ put $a_x=2|E(G[C_i])|$, and for a vertex $x\notin S$ put $a_x=0$. Let $b_x=w_x-a_x\ge0$, $A=\sum_xa_x$ and $B=\sum_xb_x$, so that $W=A+B$. Two different components $C_i,C_j$ are not joined by an edge of $G$, so every edge of $G$ that does not lie inside one of the sets $C_i$ has an endpoint outside $S$. The number $B$ counts the endpoints of such edges, and therefore
\[
B\le2\sum_{v\notin S}d_v<2nD,
\]
because every vertex outside $S$ is outside $H$.

Fix $e$ and write $A_L,b_L$ and $A_R,b_R$ for the sums of $a_x$ and $b_x$ on the two sides, so that $W_L=A_L+b_L$ and $W_R=A_R+b_R$. Then
\[
W_LW_R=A_LA_R+\bigl(A_Lb_R+A_Rb_L+b_Lb_R\bigr)\le A_LA_R+AB+B^2=A_LA_R+BW .
\]
If $A=0$ then $A_LA_R=0$ and the claim follows. Otherwise $W\ge A$ gives
\[
\frac{W_LW_R}W\le\frac{A_LA_R}{A_L+A_R}+B .
\]
Let $x$ and $y$ be the numbers of vertices of $S$ on the two sides of $e$, so $x+y=s$. A simple graph on $s_i$ vertices has fewer than $s_i^2$ edge endpoints, hence $A_L\le\sum s_i^2\le x^2$, the sum being over the components on the left, and likewise $A_R\le y^2$. The function $(p,q)\mapsto pq/(p+q)$ is increasing in each variable on $p,q\ge0$ (read as $0$ when $p+q=0$), so
\[
\frac{A_LA_R}{A_L+A_R}\le\frac{x^2y^2}{x^2+y^2}\le\frac{s^2}8 ,
\]
where the last inequality is $(x+y)^2(x^2+y^2)-8x^2y^2=(x-y)^2(x^2+4xy+y^2)\ge0$.
\end{proof}

\begin{proof}[Proof of the upper bound~\eqref{eq:upper}]
The inequality $\tau_2\le\tau_0$ is in Lemma~\ref{lem:formulas}. By Lemmas~\ref{lem:tree}, \ref{lem:cut} and~\ref{lem:components}, together with~\eqref{eq:edges} and $|E(T)|\le n-1$,
\[
\tau_0\le\Bigl(n-s+\frac sD\Bigr)\frac{s^2}8+(n-1)\cdot2nD+\frac{5Ws}{2D}.
\]
Since $s\le n$ and $W\le n^2$, the right side is at most
\[
\frac{(n-s)s^2}8+\frac{n^3}{8D}+2n^2D+\frac{5n^3}{2D}=\frac{(n-s)s^2}8+\frac{21}8\cdot\frac{n^3}D+2n^2D .
\]
On $0\le s\le n$ the function $(n-s)s^2$ has its maximum $4n^3/27$ at $s=2n/3$, so $(n-s)s^2/8\le n^3/54$. Taking $D=\sqrt n$ gives
\[
\tau_0\le\frac{n^3}{54}+\frac{21}8n^{5/2}+2n^{5/2}=\frac{n^3}{54}+\frac{37}8n^{5/2}.\qedhere
\]
\end{proof}

\begin{remark}
The maximum of $(n-s)s^2$ at $s=2n/3$ matches the barbell: two cliques on $n/3$ vertices each make up the set $S$, and the path, with about $n/3$ edges, is the tree outside it. The proof bounds every tree edge as if it were the middle edge of the path; for the barbell this costs nothing, because every path edge has the two cliques on opposite sides.
\end{remark}

\section{The lower bound}\label{sec:lower}

Let $a\ge2$ and $\ell\ge1$, and let $B(a,\ell)$ be the graph formed by two disjoint copies of $K_a$ and a path with $\ell$ edges joining a vertex $p_0$ of the first copy to a vertex $p_\ell$ of the second, through $\ell-1$ new vertices $p_1,\dots,p_{\ell-1}$. It has $2a+\ell-1$ vertices.

\begin{proposition}\label{prop:barbell}
$\displaystyle\tau_2(B(a,\ell))\ge\frac{3a(a-1)\ell+\ell^2+2}6$.
\end{proposition}

\begin{proof}
Define $f=-1$ on the first clique, $f=1$ on the second, and $f(p_j)=-1+2j/\ell$ for $0\le j\le\ell$; the two definitions agree at $p_0$ and $p_\ell$. The reflection of $B(a,\ell)$ that exchanges the cliques and maps $p_j$ to $p_{\ell-j}$ preserves degrees and changes the sign of $f$, so $\sum_vd_vf_v=0$, and~\eqref{eq:rayleigh} gives
\[
\tau_2=\frac1{\mu_2}\ge\frac{\sum_vd_vf_v^2}{\sum_{uv\in E}(f_u-f_v)^2}.
\]
Only the $\ell$ path edges contribute to the denominator, each with $(2/\ell)^2$, so the denominator equals $4/\ell$. In the numerator, each clique has total degree $a(a-1)+1$ and $f^2=1$ on it, and the $\ell-1$ inner path vertices have degree $2$, so the numerator equals
\[
2a(a-1)+2+2\sum_{j=1}^{\ell-1}\Bigl(1-\frac{2j}\ell\Bigr)^2=2a(a-1)+2+\frac{2(\ell-1)(\ell-2)}{3\ell}=2a(a-1)+\frac{2\ell}3+\frac4{3\ell},
\]
using $\sum_{j=1}^{\ell-1}(1-2j/\ell)^2=(\ell-1)-2(\ell-1)+\tfrac{2(\ell-1)(2\ell-1)}{3\ell}=\tfrac{(\ell-1)(\ell-2)}{3\ell}$. Multiplying by $\ell/4$ gives the claim.
\end{proof}

\begin{proof}[Proof of the lower bound~\eqref{eq:lower}]
For $n\ge6$, $a=\lfloor n/3\rfloor\ge2$ and $\ell=n-2a+1\ge n/3+1$, and $B_n=B(a,\ell)$ has $2a+\ell-1=n$ vertices. The first inequality in~\eqref{eq:lower} is Lemma~\ref{lem:formulas} and the second is Proposition~\ref{prop:barbell}. For the third, $a\ge(n-2)/3$ gives $a(a-1)\ge(n-2)(n-5)/9$, and $\ell\ge(n+3)/3$, so
\[
\frac{3a(a-1)\ell+\ell^2+2}6\ge\frac{a(a-1)\ell}2\ge\frac{(n-2)(n-5)(n+3)}{54}.\qedhere
\]
\end{proof}

\begin{proof}[Proof of Corollary~\ref{cor:asym}]
For $n\ge6$, Theorem~\ref{thm:main} gives
\[
\frac{(n-2)(n-5)(n+3)}{54}\le\tau_2(B_n)\le\max_G\tau_2(G)\le\max_G\tau_0(G)\le\frac{n^3}{54}+\frac{37}8n^{5/2},
\]
and $\tau_0(B_n)$ lies between the same bounds. The two outer terms are $n^3/54+O(n^{5/2})$.
\end{proof}

\section{A numerical check}\label{sec:numerics}

The script in the repository listed under Data availability computes three things. First, it computes Kemeny's constant of small barbells exactly, as $\tr Z-1$ in rational arithmetic, and compares it with~\eqref{eq:eigentime} and~\eqref{eq:resistance} evaluated in floating point; for example $\tau_0(B(3,2))=293/24$ and $\tau_0(B(6,7))=31363/222$. Second, for the barbells $B_n$ it compares the bound of Proposition~\ref{prop:barbell} with $\tau_2$ and $\tau_0$ computed from the spectrum. Table~\ref{tab:barbell} shows the ratios to $n^3$, which approach $1/54\approx0.018519$ from above. Third, on $400$ random connected graphs with $8$ to $46$ vertices (Erd\H{o}s--R\'enyi graphs, random trees, chains of cliques and lollipops with extra edges) and the thresholds $D=2,3,\sqrt n,n/4,n/2$, it checks every inequality of Section~\ref{sec:upper}: Lemma~\ref{lem:components}, the bound $B<2nD$, Lemma~\ref{lem:cut} for each edge of a breadth-first spanning tree of the quotient, Lemma~\ref{lem:tree}, the bound $R_{uv}\le\dist_T(q(u),q(v))+\delta$, and~\eqref{eq:upper}. All of them hold.

\begin{table}[ht]
\centering
\begin{tabular}{rrrccc}
\toprule
$n$ & $a$ & $\ell$ & Proposition~\ref{prop:barbell}$/n^3$ & $\tau_2(B_n)/n^3$ & $\tau_0(B_n)/n^3$\\
\midrule
30 & 10 & 11 & 0.019093 & 0.019387 & 0.021426\\
90 & 30 & 31 & 0.018718 & 0.018757 & 0.019266\\
300 & 100 & 101 & 0.018580 & 0.018583 & 0.018716\\
600 & 200 & 201 & 0.018549 & 0.018550 & 0.018614\\
\bottomrule
\end{tabular}
\caption{Balanced barbells $B_n$.}\label{tab:barbell}
\end{table}

\section*{Acknowledgements and funding}
The research presented in this paper was supported by the European Research
Council (ERC) under the European Union's Horizon Europe research and
innovation programme (grant agreement No.~101041711), by the Simons
Foundation as part of the Collaboration on the Mathematical and Scientific
Foundations of Deep Learning, by Heights Labs, by Convex Nexus Capital, by
the Israel Science Foundation (grant number 2258/19), and by the Israel
Science Foundation (ISF Grant 4101/25). Computational resources were
provided in part by Google Cloud through its education credits program
(credit 530550103).

This work was done by multiple AI agents, Claude (Anthropic) and Codex (OpenAI), connected to Math Brain~\cite{MathBrainRepo} and working together inside the Brain framework~\cite{BrainFramework}, a shared, persistent memory in which the agents record claims, proofs, failed attempts and methods, retrieve one another's results, coordinate who works on what, and audit each other's work. For more on the details behind the Brain and Math Brain, see the Math Brain paper~\cite{MathBrainPaper}, the Math Brain repository~\cite{MathBrainRepo} and the Brain framework repository~\cite{BrainFramework}.

\section*{Declaration of competing interest}
The author holds equity in Heights Labs and in Convex Nexus Capital and
has a management role at Convex Nexus Capital. These entities are included
in the funding declaration above.

\section*{Declaration of generative AI and AI-assisted technologies}
The paper comes out of Math Brain, an ongoing project on the use of AI systems for research mathematics. In September and October 2026 a GPT-family model (OpenAI), running in the Codex command-line tool, found the proof of Theorem~\ref{thm:main}; other model instances checked it. In November 2026 Claude Opus~5.5 (Anthropic), working in Claude Code, checked the proof, wrote and ran the computations of Section~\ref{sec:numerics}, and wrote the paper. The author directed the work and is responsible for the content of the paper and for decisions concerning its dissemination.

\section*{Data availability}
No data were used. The code for the computations of Section~\ref{sec:numerics}, with its output, is available at\\ \url{https://github.com/yspennstate/largest-kemeny-constant}.

The paper is part of the Math Brain project, whose public repository is \url{https://github.com/yspennstate/math-brain-public}~\cite{MathBrainRepo}.

\begin{thebibliography}{9}

\bibitem{AF}
D.~Aldous and J.~A.~Fill, \emph{Reversible Markov Chains and Random Walks on Graphs}, unfinished monograph, 2002, recompiled 2014; available at \url{https://www.stat.berkeley.edu/~aldous/RWG/book.html}.

\bibitem{BBD}
J.~Breen, S.~Butler, N.~Day, C.~DeArmond, K.~Lorenzen, H.~Qian and J.~Riesen, Computing Kemeny's constant for a barbell graph, \emph{Electron. J. Linear Algebra} \textbf{35} (2019), 583--598; doi:10.13001/1081-3810.4095.

\bibitem{BW}
G.~Brightwell and P.~Winkler, Maximum hitting time for random walks on graphs, \emph{Random Structures Algorithms} \textbf{1} (1990), 263--276.

\bibitem{CRRST}
A.~K.~Chandra, P.~Raghavan, W.~L.~Ruzzo, R.~Smolensky and P.~Tiwari, The electrical resistance of a graph captures its commute and cover times, \emph{Comput. Complexity} \textbf{6} (1996/97), 312--340.

\bibitem{FKK}
N.~Faught, M.~Kempton and A.~Knudson, A 1-separation formula for the graph Kemeny constant and Braess edges, \emph{J. Math. Chem.}, doi:10.1007/s10910-021-01294-8; arXiv:2108.01061.

\bibitem{Fe}
U.~Feige, A tight upper bound on the cover time for random walks on graphs, \emph{Random Structures Algorithms} \textbf{6} (1995), 51--54.

\bibitem{KS}
J.~G.~Kemeny and J.~L.~Snell, \emph{Finite Markov Chains}, Van Nostrand, Princeton, 1960.

\bibitem{MathBrainRepo} Y.~Shmalo, \emph{Math Brain: public snapshot, code and tests}, 2026, \url{https://github.com/yspennstate/math-brain-public}.

\bibitem{MathBrainPaper} Y.~Shmalo, \emph{Math Brain: shared, auditable memory for multi-agent mathematical research}, October 2026, \url{https://github.com/yspennstate/math-brain-public/tree/main/papers/math-brain-2026-10-24}.

\bibitem{BrainFramework} Y.~Shmalo, \emph{Brain Framework}, public software repository, 2026, \url{https://github.com/yspennstate/brain-framework-public}.

\end{thebibliography}
\end{document}
